\documentclass[11pt]{article}
\usepackage[a4paper,margin=25mm]{geometry}
\usepackage[T1]{fontenc}
\usepackage{lmodern}
\usepackage{booktabs}
\usepackage{graphicx}
\usepackage{hyperref}
\usepackage{xcolor}
\usepackage{listings}
\lstset{basicstyle=\ttfamily\small, columns=fullflexible, keepspaces=true,
  breaklines=true, frame=single, framerule=0.3pt, rulecolor=\color{gray}}
\newcommand{\code}[1]{\texttt{#1}}

\title{\code{chanmap.py}: the array on the face of the crystal}
\author{BGOTSCAN}
\date{2026-10-15}

\begin{document}
\maketitle

\section{Purpose}

One figure of the 64 SiPMs at their positions on the face of the crystal,
from \code{src/channel\_map.csv}, each cell labelled with its digitizer
channel, the twelve dead channels in grey and the four high-gain channels
in orange. It keeps the dead block in sight in the notes.

\section{Usage}

\begin{lstlisting}
python3 src/py/chanmap.py [--dead LIST] [--ringers LIST] [--save NAME]
\end{lstlisting}

\begin{center}
\begin{tabular}{l l p{80mm}}
\toprule
option & default & meaning \\
\midrule
\code{--dead LIST} & 19,20,21,22,23,25,27,33,35,41,46,47 & dead channels (\code{chanstat\_doc.tex}) \\
\code{--ringers LIST} & 11,12,48,49 & high-gain channels (\code{spe\_doc.tex}) \\
\code{--save NAME} & & save the figure as \code{src/py/plots/NAME\_chanmap.png} \\
\bottomrule
\end{tabular}
\end{center}

Each SiPM is drawn as a 6~mm square at its $(x, y)$; the axes are in mm,
$x$ along the columns, $y$ decreasing with the row, as the map file
defines them.

\section{Example}

\begin{lstlisting}
python3 src/py/chanmap.py --save array      # docs/figs/chanmap.png
\end{lstlisting}

\end{document}
